\section{Warum dieses Dokument eingeschoben ist}

A080 hat die Einheiten als Buchhaltung behandelt: Brückenkonstanten, die
Trennung von Verhältnis und Absolutwert. Das ist die \emph{operative} Seite.
Dieses Dokument liefert die \emph{fundamentale} Seite, die A080 voraussetzt,
aber nicht ausführt: warum in dieser Theorie überhaupt ein natürliches
Einheitensystem existiert, in dem alle dimensionsbehafteten Konstanten den Wert
Eins annehmen.

Es steht vor dem Gravitationsdokument (A145) und vor der Quantenmechanik (A160),
weil beide es benutzen: G wird aus $\xi$ und einer Masse gewonnen, und das ist
nur in natürlichen Einheiten eine Aussage über Struktur statt über
Zahlenkonventionen.

\section{Der Ursprung liegt in der Dualität}

\textbf{[K]} Das natürliche Einheitensystem der Theorie ist keine Wahl, sondern
eine Folge der ersten Setzung. Aus $\tilde T\cdot m = 1$ (A020) folgt, dass
Zeit und Masse reziprok sind -- dieselbe Größe in zwei Lesarten. Damit ist die
Umrechnung zwischen ihnen keine physikalische Konstante, sondern die Identität.

\textbf{[B]} In der üblichen Physik leisten drei Konstanten die Umrechnung
zwischen vier Grunddimensionen:
\[
  c:\ \text{Länge}\leftrightarrow\text{Zeit},\quad
  \hbar:\ \text{Masse}\leftrightarrow\text{Zeit},\quad
  k_B:\ \text{Temperatur}\leftrightarrow\text{Energie}.
\]
Sobald aber Zeit und Masse dieselbe Größe sind, kollabiert das System: es gibt
nur noch eine unabhängige Dimension, und diese drei Umrechnungsfaktoren werden
zu Eins.

\textbf{[SETZUNG]} Ein Wort zur Abgrenzung, das hierher gehört: dieser
Mechanismus betrifft die \emph{dimensionsbehafteten} Brückenkonstanten
$c, \hbar, k_B$ (und, über die Masse-Länge-Verbindung, $G$; A145). Die
Feinstrukturkonstante $\alpha$ ist \emph{nicht} dabei. Sie ist dimensionslos,
und ihre Reduktion auf Eins hat einen anderen Grund -- die elektromagnetische
Dualität, nicht die Zeit-Masse-Dualität. Das ist ein eigenes Kapitel und wird in
A130 geführt, nicht hier. Dieses Dokument behauptet über $\alpha$ nichts; es
sagt nur, warum $c, \hbar, k_B$ zu Eins werden.

\textbf{[K]} In natürlichen Einheiten gilt daher
\[
  [M] = [E] = [T^{-1}] = [L^{-1}],
\]
Masse, Energie, inverse Zeit und inverse Länge sind eine einzige Dimension. Das
ist nicht vereinfachend gesetzt, sondern die direkte Konsequenz der Dualität.

\section{Warum das mehr ist als eine Konvention}

\textbf{[B]} Man kann in jeder Theorie $\hbar = c = 1$ setzen; das allein ist
folgenlos. Der Unterschied hier ist, dass die \emph{Reduktion auf eine einzige
Dimension} eine inhaltliche Behauptung ist und nicht eine Schreibweise.

In der Standardphysik bleibt nach $\hbar = c = 1$ noch eine Massendimension
übrig -- Energie ist eine echte Einheit, die man wählen muss. Hier bleibt
\emph{nichts} zu wählen: die eine verbleibende Dimension ist durch $\xi$
skaliert (A020), und die absolute Skala ist die einzige verbliebene Eingabe
(A080, A130).

\textbf{[SETZUNG]} Damit ist die Buchhaltung aus A080 an ihren Grund gebunden:
die Brückenkonstanten sind nicht deshalb gleich Eins, weil man sie so setzt,
sondern weil die Dualität ihre Dimensionen identifiziert. Die Wahl $=1$ ist
das Ablesen einer Identität, nicht eine Normierung.

\section{Die vier Kleidungen, fundamental gelesen}

\textbf{[K]} A080 hat die Identität
\[
  m_0 c^2 = E_0 = k_B T_0 = \hbar/\tilde T
\]
als exakte Umrechnung eingeführt. Hier ist ihr Grund: die vier Ausdrücke sind
nicht vier Größen, die zufällig gleich sind, sondern eine Größe in vier
Kleidungen, weil die Dualität Masse, Energie, Temperatur und inverse Zeit auf
dieselbe Dimension legt.

Die Umrechnungsfaktoren $c^2$, $k_B$, $\hbar$ sind die Nähte zwischen den
Kleidungen. In natürlichen Einheiten verschwinden die Nähte, weil es nur ein
Kleidungsstück gibt.

\section{Was daraus für die folgenden Dokumente folgt}

\textbf{[B]} Drei Konsequenzen, die A145 und A160 voraussetzen:

Erstens ist jede Masse eine inverse Länge und umgekehrt. Das erlaubt es, eine
Gravitationskopplung, die Massen mit Längen verbindet, allein durch $\xi$
auszudrücken -- der Inhalt von A145.

Zweitens ist jede dimensionsbehaftete Konstante entweder Eins oder eine
Kalibrierung. Es gibt keinen dritten Fall. Wo im Korpus eine solche Konstante
einen anderen Wert trägt, ist das ein SI-Artefakt, kein Naturfaktum.

Drittens sind Verhältnisse gegenüber der Skalenwahl invariant (A080), und die
Skalenwahl ist der \emph{einzige} verbleibende Freiheitsgrad des
Einheitensystems. Alles Dimensionslose ist damit skalenfrei bestimmt oder gar
nicht.

\section{Prüfskript}

\begin{itemize}
  \item Dass $\tilde T\cdot m=1$ das Vier-Dimensionen-System auf eine
    Dimension reduziert, gezeigt an der Dimensionsbuchhaltung; und dass die
    Planck-Größen genau die Umrechnungsfaktoren des reduzierten Systems sind:
    \texttt{python/A\_Serie\_Skripte/a085\_natuerliche\_einheiten.py}
\end{itemize}

\section{Was offen bleibt}

\textbf{Erstens bleibt die absolute Skala eine Eingabe.} Die Reduktion auf eine
Dimension ist hergeleitet; welche Größe diese eine Dimension hat, ist es nicht.
Das ist dieselbe Lücke wie in A020 (Größenordnung von $\xi$) und A130
(Kalibrierung von $E_0$), hier in ihrer allgemeinsten Form.

\textbf{Zweitens ist die Identifikation von Temperatur mit Energie eine
thermodynamische Zusatzannahme.} $k_B=1$ setzt voraus, dass Temperatur eine
Energie pro Freiheitsgrad ist; das ist Standard, aber es folgt nicht aus der
Dualität, sondern tritt neben sie.

\section{Ersetzt}

Dieses Dokument ergänzt A080 um die fundamentale Begründung und nimmt die
entsprechenden Teile auf von: Dok.~013 (SI), Dok.~014 (natürlich/SI),
Dok.~015 (Systematik der Natureinheiten), Dok.~012 (Einheitenteil).
Eingearbeitet: P40.

\section{Verwendung von KI-Werkzeugen}

Dieses Dokument ist unter Verwendung KI-gestützter Werkzeuge entstanden. Die
Fragestellungen, die Setzungen und alle inhaltlichen Entscheidungen stammen vom
Autor; die Werkzeuge formulieren aus, rechnen nach und erstellen Prüfskripte.
Die Verantwortung für Inhalt und Fehler liegt vollständig beim Autor. Der
ausführliche Wortlaut steht in A010.
